\chapter{Designing Filters}

Up to now we have mostly \emph{analyzed} signals: we took something the world
handed us --- a chord, an accelerometer trace, an image --- and asked what
frequencies were hiding inside it. A filter does the opposite. A filter is a
machine we \emph{design on purpose} to keep some frequencies and throw others
away. It is the first piece of signal processing where you are no longer a
spectator; you are an engineer with a target in mind.

And you already own all the parts. In the chapter on LTI systems we learned that
any linear, time-invariant box is completely described by its impulse response
$h[n]$, that running a signal through it is convolution $y = x \conv h$, and ---
the magic key --- that convolution in time is \emph{multiplication} in frequency:
\[
Y(\omega) \;=\; X(\omega)\, H(\omega).
\]
That single equation \emph{is} the theory of filtering. The factor $H(\omega)$,
the frequency response, is a knob you get to draw: a curve that says ``multiply
this frequency by $1$, that one by $0$, this one by $0.5$.'' Designing a filter
means choosing an $H(\omega)$ shaped the way you want, and then finding a real,
finite, runnable $h[n]$ (or a difference equation) that produces it. This chapter
is about how to do that choosing well.

\begin{focusbox}
A filter multiplies each frequency of the input by a number you chose in advance.
The wish-list curve is $H(\omega)$. ``Filter design'' is the craft of turning that
wish into a short list of coefficients a chip can actually multiply and add ---
trading away a little of the wish to make it cheap, stable, and causal.
\end{focusbox}

\section{What a filter is supposed to do}
\label{ch16:what}

Forget formulas for a second and think in pictures. A signal's spectrum is a
landscape: energy piled up at some frequencies, flat at others. A filter is a
\emph{mask} you lay over that landscape. Where the mask is $1$ the energy passes
untouched; where the mask is $0$ the energy is erased; in between it is scaled.
Four mask shapes cover almost everything you will ever want.

\begin{itemize}
\item \textbf{Low-pass}: keep low frequencies, kill high ones. The mask is a
plateau near $\omega=0$ that drops to zero. This is the workhorse --- smoothing,
anti-aliasing, removing hiss, taming sensor jitter.
\item \textbf{High-pass}: keep high frequencies, kill low ones. The mask is zero
near DC and rises to a plateau. Use it to remove a slow drift or a DC offset
while keeping the wiggles.
\item \textbf{Band-pass}: keep a middle band, kill everything above and below. A
plateau floating between two cliffs. Radios tune with these; so does an EMG front
end that only wants the $20$--$450$\,Hz muscle band.
\item \textbf{Band-stop} (notch): the inverse --- kill one narrow band, keep the
rest. The classic job is murdering $50/60$\,Hz mains hum without touching the
signal around it.
\end{itemize}

\begin{figure}[H]
\centering
\begin{tikzpicture}[scale=1.0]
\begin{axis}[
  width=13.5cm, height=4.2cm,
  xmin=0, xmax=1, ymin=-0.1, ymax=1.25,
  xtick={0,0.5,1}, xticklabels={$0$,$\pi/2$,$\pi$},
  ytick={0,1},
  xlabel={frequency $\omega$}, ylabel={$|H|$},
  every axis plot/.append style={thick},
  axis lines=box, enlargelimits=false,
  title={\small Four ideal masks (low/high/band-pass, band-stop)},
]
% low-pass
\addplot[sigblue, samples=2, domain=0:0.3] {1};
\addplot[sigblue, samples=2, domain=0.3:1] {0};
% high-pass
\addplot[specred, samples=2, domain=0:0.7] {0};
\addplot[specred, samples=2, domain=0.7:1] {0.92};
% band-pass
\addplot[quotientorange, samples=2, domain=0:0.42] {0};
\addplot[quotientorange, samples=2, domain=0.42:0.58] {0.84};
\addplot[quotientorange, samples=2, domain=0.58:1] {0};
\end{axis}
\end{tikzpicture}
\caption{Idealized magnitude responses. Blue: low-pass. Red: high-pass (drawn
slightly low so the lines do not overlap). Orange: band-pass. A band-stop is just
``$1$ minus a band-pass''.}
\label{ch16:fig-masks}
\end{figure}

\begin{intuitionbox}
A high-pass is a low-pass turned inside out, and a band-stop is a band-pass turned
inside out. Spectrally, ``keep the rest'' means subtract from the all-pass mask
$H=1$: if $H_{\text{LP}}(\omega)$ keeps the lows, then $1 - H_{\text{LP}}(\omega)$
keeps the highs. In time that is $\delta[n] - h_{\text{LP}}[n]$. So if you can
build one good low-pass, you can build all four. We will lean on this constantly.
\end{intuitionbox}

\begin{connectionbox}
The masking picture is the convolution theorem from the Fourier chapter wearing
work clothes. Multiplying spectra ($Y = XH$) corresponds to convolving in time
($y = x \conv h$). Everything in this chapter is just: ``what time-domain $h[n]$
multiplies the spectrum by the shape I want?'' That is an inverse-Fourier
question, and the linear-algebra book's view of the Fourier transform as an
orthogonal change of basis is exactly why the answer is unique.
\end{connectionbox}

\subsection{Why the perfect mask is impossible}
\label{ch16:why-impossible}

The masks in Figure~\ref{ch16:fig-masks} have vertical cliffs: pass at full
strength up to a frequency, then \emph{instantly} zero. That perfect ``brick
wall'' is the dream. Here is the hard truth that organizes the whole field: you
cannot build it with a finite, real filter, and you cannot even get arbitrarily
close without paying.

The reason is the uncertainty relationship between time and frequency. A brick
wall in frequency is a $\rect$ function. Its inverse Fourier transform is a
$\sinc$ --- and the $\sinc$ is infinitely long and two-sided.

\begin{whybox}
An ideal low-pass with cutoff $\omega_c$ has frequency response
$H(\omega) = \rect(\omega/2\omega_c)$ (one inside the band, zero outside).
Inverse-transforming gives the impulse response
\[
h[n] \;=\; \frac{\omega_c}{\pi}\,\sinc\!\left(\frac{\omega_c\, n}{\pi}\right)
       \;=\; \frac{\sin(\omega_c n)}{\pi n}.
\]
Look at what this demands. It is nonzero for \emph{all} $n$, from $-\infty$ to
$+\infty$. To compute even one output sample you would need every past
\emph{and every future} input sample, weighted by a slowly decaying ripple. That
filter is non-causal (needs the future) and infinitely long (needs infinite
memory and infinite multiplies). Real hardware has neither. Every practical
design is a negotiation: how do we chop, shift, and smooth that ideal $\sinc$
into something finite and causal while keeping the mask close to what we wanted?
\end{whybox}

\begin{figure}[H]
\centering
\begin{tikzpicture}
\begin{axis}[
  width=13.5cm, height=4.6cm,
  xmin=-12, xmax=12, ymin=-0.09, ymax=0.34,
  xtick={-10,-5,0,5,10}, ytick={0,0.3},
  xlabel={sample index $n$}, ylabel={$h[n]$},
  axis lines=box, enlargelimits=false,
  title={\small Ideal low-pass impulse response: the sinc never ends},
]
\addplot[sigblue, only marks, mark=*, mark size=1.2pt, samples at={-12,...,12}]
  ({x}, {x==0 ? 0.3 : sin(deg(0.3*pi*x))/(pi*x)});
\addplot[sigblue, ycomb, samples at={-12,...,12}]
  ({x}, {x==0 ? 0.3 : sin(deg(0.3*pi*x))/(pi*x)});
\draw[gray, dashed] (axis cs:-12,0)--(axis cs:12,0);
\end{axis}
\end{tikzpicture}
\caption{The impulse response of the ideal low-pass ($\omega_c = 0.3\pi$) is a
sampled sinc. It is symmetric about $n=0$, decays only as $1/n$, and goes on
forever in both directions. Truncating and shifting it is the entire ``window
method'' below.}
\label{ch16:fig-ideal-sinc}
\end{figure}

\section{Two families: FIR and IIR}
\label{ch16:fir-iir}

Every digital filter falls into one of two families, distinguished by one
question: \emph{does the impulse response ever truly end?}

A \textbf{Finite Impulse Response} (FIR) filter has an $h[n]$ that is nonzero for
only finitely many taps. You compute each output as a plain weighted sum of recent
inputs:
\[
y[n] \;=\; \sum_{k=0}^{M} b_k\, x[n-k].
\]
There is no feedback. Kick it with an impulse and the output rings for exactly
$M+1$ samples, then is dead silent forever. The coefficients $b_k$ \emph{are} the
impulse response: $h[k] = b_k$.

An \textbf{Infinite Impulse Response} (IIR) filter feeds its own past outputs
back in:
\[
y[n] \;=\; \sum_{k=0}^{M} b_k\, x[n-k] \;-\; \sum_{k=1}^{P} a_k\, y[n-k].
\]
Because the output depends on earlier outputs, a single impulse can echo around
the feedback loop forever (decaying, if the filter is stable). A two-tap IIR can
produce an impulse response that never exactly reaches zero --- hence
``infinite''.

\begin{intuitionbox}
FIR is a \emph{moving window} that only ever looks at the input: a tapped delay
line you weight and sum. IIR is a \emph{resonator}: it has internal state that
keeps ringing, like a struck bell or an RC circuit relaxing. The bell's tone
outlives the strike --- that is feedback, that is the ``infinite'' tail.
\end{intuitionbox}

\begin{happeningbox}
Both filters are still just convolution, $y = x \conv h$. The difference is
\emph{how} you compute that convolution. FIR computes it directly, multiplying out
a finite list of taps. IIR computes the \emph{same kind of} convolution
implicitly through a recurrence: a short feedback formula that, if you unrolled it
forever, would reproduce an infinitely long $h[n]$. IIR is a clever shorthand ---
it buys a long, sharp impulse response with very few coefficients, and the bill
for that bargain is the risk of instability and nonlinear phase.
\end{happeningbox}

\subsection{The tradeoffs, plainly}
\label{ch16:tradeoffs}

Neither family is ``better''; they trade different things. Here is the
ledger.

\begin{center}
\renewcommand{\arraystretch}{1.3}
\begin{tabular}{p{4.1cm} p{4.5cm} p{4.5cm}}
\toprule
 & \textbf{FIR} & \textbf{IIR} \\
\midrule
Sharpness per coefficient & Low --- needs many taps for a steep cutoff
 & High --- a few coefficients give a very steep cutoff \\
Stability & Always stable (no feedback) & Can blow up if poorly designed or
 quantized \\
Phase & Can be \emph{exactly} linear (constant delay) & Nonlinear phase; warps
 waveform shape \\
Memory / compute & More (a long delay line) & Less \\
Design intuition & Direct: shape the taps & Indirect: place poles and zeros \\
Analog heritage & None really & Mirrors classic analog filters (RC, LC,
 Butterworth) \\
\bottomrule
\end{tabular}
\end{center}

\begin{trapbox}
The seductive trap is to reach for IIR \emph{because} it is cheap, then discover
your nicely-shaped signal has its peaks shifted relative to its valleys. IIR
filters delay different frequencies by different amounts (nonlinear phase). For
audio loudness that is usually fine --- the ear is fairly deaf to phase. For
anything where the \emph{shape} of the waveform matters --- an ECG QRS complex, a
gesture's sharp accelerometer spike, lining up two channels --- that smearing is a
real bug. When shape matters, prefer FIR.
\end{trapbox}

\begin{goalbox}
For the wrist gesture controller these tradeoffs are not academic. We will run a
low-pass on the accelerometer to separate gravity (slow tilt) from motion (fast),
a band-pass on any EMG/FMG channel to isolate the muscle band, and a notch to kill
mains hum. On a small microcontroller, compute and memory are scarce, which pulls
toward IIR. But a \emph{tap} or \emph{flick} gesture is recognized by the precise
timing and shape of an acceleration spike, and IIR's phase smearing can blur two
near-simultaneous taps into one. So a typical wrist build mixes both: cheap IIR
biquads for the EMG band-pass and the hum notch (where only the envelope matters),
and a short linear-phase FIR on the accelerometer (where spike timing must stay
aligned across axes). Choosing per channel, with the table above in hand, is the
real engineering.
\end{goalbox}

\section{Specifying a filter before you design it}
\label{ch16:spec}

You cannot hit a target you have not drawn. Before any design method, you write a
\emph{specification}: a tolerance template the magnitude response must live inside.
Because the brick wall is impossible (Section~\ref{ch16:why-impossible}), the spec
deliberately leaves slack. Five numbers describe a low-pass spec.

\begin{itemize}
\item \textbf{Passband edge} $\omega_p$: up to here, keep the signal.
\item \textbf{Stopband edge} $\omega_s$: beyond here, kill the signal.
\item \textbf{Transition band} $(\omega_p, \omega_s)$: the ramp in between, where
we make no promises. Its \emph{width} is what costs you taps. Narrow transition =
expensive filter.
\item \textbf{Passband ripple} $\delta_p$: how much the gain is allowed to wobble
around $1$ in the passband (often quoted in dB).
\item \textbf{Stopband attenuation} $\delta_s$: how far down the leftover must be
pushed in the stopband (e.g. $-60$\,dB means leakage no larger than $0.001$).
\end{itemize}

\begin{figure}[H]
\centering
\begin{tikzpicture}
\begin{axis}[
  width=13.5cm, height=5.4cm,
  xmin=0, xmax=1, ymin=-0.08, ymax=1.22,
  xtick={0.32,0.5}, xticklabels={$\omega_p$,$\omega_s$},
  ytick={0,0.06,0.94,1,1.06}, yticklabels={$0$,$\delta_s$,$1-\delta_p$,$1$,$1+\delta_p$},
  xlabel={frequency $\omega$}, ylabel={$|H(\omega)|$},
  axis lines=box, enlargelimits=false,
  title={\small A low-pass tolerance template},
]
% passband tolerance band
\addplot[draw=none, fill=sigblue!18, domain=0:0.32, samples=2] {1.06} \closedcycle;
\addplot[draw=none, fill=white, domain=0:0.32, samples=2] {0.94} \closedcycle;
% stopband tolerance
\addplot[draw=none, fill=specred!18, domain=0.5:1, samples=2] {0.06} \closedcycle;
% transition shading
\addplot[draw=none, fill=gray!12, domain=0.32:0.5, samples=2] {1.22} \closedcycle;
% a plausible response curve
\addplot[sigblue, very thick, samples=120, domain=0:1]
  {1.0/sqrt(1+ (x/0.40)^12) + 0.018*cos(deg(40*x))*exp(-3*x)};
% edge guides
\draw[gray,dashed] (axis cs:0.32,-0.08)--(axis cs:0.32,1.22);
\draw[gray,dashed] (axis cs:0.5,-0.08)--(axis cs:0.5,1.22);
\node[gray] at (axis cs:0.41,1.13) {\scriptsize transition};
\end{axis}
\end{tikzpicture}
\caption{The spec is a corridor, not a curve. The response (blue) must stay inside
the passband ripple band (top left), may do anything in the gray transition, and
must stay under the stopband ceiling $\delta_s$ (bottom right). Tightening any of
$\delta_p$, $\delta_s$, or the transition width raises the filter's order.}
\label{ch16:fig-spec}
\end{figure}

\begin{formalbox}
Magnitudes are usually quoted in decibels, $|H|_{\text{dB}} = 20\log_{10}|H|$.
Handy landmarks: a gain of $1$ is $0$\,dB; the half-power point $|H| = 1/\sqrt2$ is
$-3$\,dB (this defines the conventional ``cutoff'' for analog-style filters);
$|H| = 0.1$ is $-20$\,dB; $|H| = 0.001$ is $-60$\,dB. Passband ripple of
$0.1$\,dB means the gain stays within $\pm 1.2\%$ of unity. The decibel scale is
used because attenuation requirements span many orders of magnitude and our
perception of ``how much is left'' is roughly logarithmic.
\end{formalbox}

\section{Linear phase: the FIR superpower}
\label{ch16:linphase}

Before designing anything, understand the one property that makes FIR worth its
extra taps: it can have \emph{exactly linear phase}, which means \emph{constant
delay for every frequency} --- and that means it does not distort waveform shape.

What is phase doing? The frequency response is complex,
$H(\omega) = |H(\omega)|\,e^{\jj\phi(\omega)}$. The magnitude $|H|$ says how much
each frequency is scaled; the phase $\phi(\omega)$ says how much each frequency is
\emph{delayed}. A sinusoid at frequency $\omega$ passing through the filter comes
out shifted by $\phi(\omega)$ radians, i.e. delayed in time by
$-\phi(\omega)/\omega$ samples. This is the \emph{group delay}.

\begin{whybox}
If $\phi(\omega) = -\alpha\,\omega$ is a straight line through the origin, then the
delay $-\phi(\omega)/\omega = \alpha$ is the \emph{same constant} for every
frequency. Every component of the signal is held back by exactly $\alpha$ samples,
so they all stay in step. The waveform comes out delayed but otherwise
\emph{unchanged in shape}. If instead $\phi(\omega)$ bends, low and high
components are delayed by different amounts, they slide out of alignment, and a
crisp spike spreads into a lopsided smear --- even when the magnitude response is
perfect. Constant delay = no shape distortion. That is the prize.
\end{whybox}

Here is the key fact: an FIR filter has exactly linear phase \emph{whenever its taps
are symmetric} (or antisymmetric) about their center, $h[k] = h[M-k]$. Symmetry is
a geometric condition you can guarantee by construction. The ideal sinc in
Figure~\ref{ch16:fig-ideal-sinc} is symmetric about $n=0$; if we keep that
symmetry while truncating and just shift the whole thing right so it starts at
$n=0$ (making it causal), linear phase is preserved automatically. The shift is
exactly the constant group delay $\alpha = M/2$ samples.

\begin{intuitionbox}
Symmetric taps mean the filter treats ``a bit before the center'' and ``a bit
after the center'' identically. There is no front-to-back lopsidedness to make one
frequency lag another. The filter's only effect on timing is a bulk delay equal to
half its length. IIR filters, built from feedback, cannot generally be symmetric
in this way, which is the deep reason they have nonlinear phase.
\end{intuitionbox}

\begin{connectionbox}
This is the inner-product / symmetry idea from the linear algebra book. A
real-symmetric tap vector projects onto cosines and sines in a balanced way; the
imaginary parts of the response arrange themselves into a pure linear phase term
$e^{-\jj\alpha\omega}$ that factors out, leaving a purely real ``zero-phase''
gain. Symmetry in the coefficient vector becomes linearity in the phase.
\end{connectionbox}

\section{FIR design by the window method}
\label{ch16:window}

Now we actually build something. The window method is the most intuitive FIR
design technique, and it falls right out of everything above:

\begin{enumerate}
\item Start from the ideal filter's impulse response (the infinite sinc).
\item \textbf{Truncate} it to $M+1$ taps --- chop off the tails.
\item \textbf{Shift} it to start at $n=0$ so the filter is causal.
\item Multiply the kept taps by a smooth \textbf{window} to tame the damage done by
the chop.
\end{enumerate}

Steps 1--3 alone already give a working filter. Let us see why step 4 is needed.

\subsection{Truncation, and the ripples it causes}
\label{ch16:gibbs}

Chopping the sinc to a finite length is, in the frequency domain, multiplying the
ideal $\rect$ mask by something not nice at all. Truncating $h[n]$ to a window of
length $N$ is the same as \emph{multiplying} the infinite ideal $h[n]$ by a
rectangular window $w[n]$ (one inside the kept range, zero outside). By the
convolution theorem, multiplication in time is convolution in frequency:
\[
H_{\text{trunc}}(\omega) \;=\; H_{\text{ideal}}(\omega) \conv W(\omega).
\]
The spectrum of a rectangular window is a $\sinc$-shaped curve with a tall main
lobe and slowly-decaying side lobes. Convolving the clean brick wall with that
smears the cliff into a ramp and --- worse --- sprinkles \emph{ringing ripples} on
both sides of the edge. This overshoot near the cliff is the \textbf{Gibbs
phenomenon}, and it is stubborn: making the filter longer makes the ripples
narrower but \emph{not} shorter. The first overshoot stays near $9\%$ no matter
how many taps you add.

\begin{figure}[H]
\centering
\begin{tikzpicture}
\begin{axis}[
  width=13.5cm, height=5.2cm,
  xmin=0, xmax=1, ymin=-0.12, ymax=1.2,
  xtick={0,0.3,0.5,1}, xticklabels={$0$,$\omega_c$,,$\pi$},
  ytick={0,1},
  xlabel={frequency $\omega$}, ylabel={$|H(\omega)|$},
  axis lines=box, enlargelimits=false,
  legend style={font=\scriptsize, at={(0.98,0.95)}, anchor=north east},
  title={\small Truncation rings (Gibbs); a window smooths it},
]
% rectangular-window result: ringing approximation of a step at 0.3
\addplot[specred, thick, samples=180, domain=0.001:1]
  {0.5 - 0.62*( (x-0.3) ) - 0.16*sin(deg(60*(x-0.3)))/(60*(x-0.3)+0.0001) };
\addplot[sigblue, very thick, samples=120, domain=0:1]
  {1/sqrt(1+(x/0.315)^16)};
\legend{rectangular (rings), windowed (smooth)}
\draw[gray,dashed] (axis cs:0.3,-0.12)--(axis cs:0.3,1.2);
\end{axis}
\end{tikzpicture}
\caption{Schematic of the window method's effect near the cutoff. A bare
rectangular truncation (red) overshoots and rings around the edge. A tapered
window (blue) trades a slightly wider transition for a much cleaner, ripple-free
stopband. The plotted curves are illustrative shapes, not an exact computation.}
\label{ch16:fig-gibbs}
\end{figure}

\subsection{Why a window helps, and which one}
\label{ch16:windows}

The cure is to stop chopping abruptly. Instead of multiplying by a hard
rectangle, multiply the ideal taps by a window $w[n]$ that \emph{tapers smoothly}
to zero at the ends. A gentler taper has a spectrum with much lower side lobes, so
the smearing it convolves into the response has far weaker ripples.

\begin{whybox}
Everything hinges on one tradeoff inside the window's own spectrum: \emph{main-lobe
width versus side-lobe height}. A wide, smooth taper (low side lobes) gives a
clean stopband --- deep attenuation, tiny ripples --- but its fat main lobe
smears the cliff into a wide, gentle transition band. A narrow taper (the
rectangle is the extreme) gives the sharpest possible transition but tall side
lobes and ugly ringing. \emph{You are buying stopband cleanliness with transition
width.} Choosing a window is choosing where on that line to sit.
\end{whybox}

A few standard windows, from sharpest-but-leakiest to smoothest:

\begin{center}
\renewcommand{\arraystretch}{1.25}
\begin{tabular}{p{2.8cm} p{3.2cm} p{6.2cm}}
\toprule
\textbf{Window} & \textbf{Peak side lobe} & \textbf{Character} \\
\midrule
Rectangular & $-13$\,dB & sharpest transition, worst stopband (rings) \\
Hann & $-31$\,dB & gentle all-rounder, easy default \\
Hamming & $-41$\,dB & like Hann but flatter stopband near the edge \\
Blackman & $-58$\,dB & very clean stopband, wider transition \\
Kaiser ($\beta$) & tunable & one knob $\beta$ slides the whole tradeoff \\
\bottomrule
\end{tabular}
\end{center}

The Kaiser window deserves a note: its single parameter $\beta$ lets you dial
exactly between ``sharp'' and ``clean'', and Kaiser gave simple formulas that take
your spec (desired stopband attenuation and transition width) and spit out both
$\beta$ \emph{and} the number of taps you need. That makes it the practical
default when you have a real spec to hit.

\begin{figure}[H]
\centering
\begin{tikzpicture}
\begin{axis}[
  width=13.5cm, height=4.4cm,
  xmin=0, xmax=1, ymin=-0.05, ymax=1.1,
  xtick={0,0.5,1}, xticklabels={$0$,$M/2$,$M$},
  ytick={0,1},
  xlabel={tap index $n$ (normalized)}, ylabel={$w[n]$},
  axis lines=box, enlargelimits=false,
  legend style={font=\scriptsize, at={(0.5,-0.32)}, anchor=north, legend columns=4},
  title={\small Window shapes: how each taper falls to zero},
]
\addplot[gray, thick, samples=2, domain=0:1] {1}; % rectangular
\addplot[sigblue, thick, samples=80, domain=0:1] {0.5-0.5*cos(deg(2*pi*x))}; % Hann
\addplot[specred, thick, samples=80, domain=0:1] {0.54-0.46*cos(deg(2*pi*x))}; % Hamming
\addplot[quotientorange, thick, samples=80, domain=0:1]
  {0.42-0.5*cos(deg(2*pi*x))+0.08*cos(deg(4*pi*x))}; % Blackman
\legend{rect, Hann, Hamming, Blackman}
\end{axis}
\end{tikzpicture}
\caption{The same length, four tapers. The rectangle keeps every tap at full
weight then drops to zero at the ends --- that cliff is what rings. Hann, Hamming,
and Blackman fade the end taps out, each a little more gently, buying lower side
lobes at the cost of a wider main lobe.}
\label{ch16:fig-windows}
\end{figure}

\subsection{Worked design: a low-pass with the Hamming window}
\label{ch16:worked-fir}

Let us design a real filter end to end. \textbf{Spec}: sampling rate
$f_s = 1000$\,Hz; pass everything below about $100$\,Hz; suppress everything above
about $200$\,Hz by at least $40$\,dB. (This is exactly the kind of accelerometer
low-pass we will want.)

\textbf{Step 1 --- normalize.} Digital filters live in normalized frequency
$\omega = 2\pi f / f_s$ (radians/sample), or in cycles/sample $f/f_s$. The natural
cutoff to aim the sinc at is the \emph{middle} of the transition band,
$f_c = 150$\,Hz, so $\omega_c = 2\pi (150/1000) = 0.3\pi$.

\textbf{Step 2 --- choose length from the window.} A Hamming window gives about
$53$\,dB stopband attenuation (plenty for our $40$\,dB) and a transition width of
roughly $\Delta\omega \approx 6.6\pi/M$ in normalized terms. We want the
transition to fit inside $100$--$200$\,Hz, i.e.
$\Delta\omega = 2\pi(100/1000) = 0.2\pi$. Solving $0.2\pi = 6.6\pi/M$ gives
$M \approx 33$, so take $M=34$ (an even $M$ gives an odd number of taps $M+1=35$,
which centers a tap exactly at the middle --- convenient for a clean integer
delay).

\textbf{Step 3 --- the ideal taps, truncated and shifted.} Center the sinc at
$\alpha = M/2 = 17$ and evaluate:
\[
h_{\text{ideal}}[n] \;=\; \frac{\sin\!\big(\omega_c (n-\alpha)\big)}{\pi (n-\alpha)},
\qquad n = 0, 1, \dots, 34,
\]
with the removable point $h_{\text{ideal}}[\alpha] = \omega_c/\pi$.

\textbf{Step 4 --- apply the window.} Multiply tap by tap:
\[
h[n] \;=\; h_{\text{ideal}}[n]\cdot w[n], \qquad
w[n] = 0.54 - 0.46\cos\!\Big(\frac{2\pi n}{M}\Big).
\]
That is the filter. Thirty-five symmetric coefficients, exactly linear phase
(constant delay of $17$ samples = $17$\,ms here), guaranteed stable. In code:

\begin{lstlisting}[language=Python]
import numpy as np

fs   = 1000.0
fc   = 150.0                 # cutoff at middle of transition band
M    = 34                    # filter order -> M+1 = 35 taps
wc   = 2*np.pi*fc/fs         # = 0.3*pi rad/sample
n    = np.arange(M+1)
a    = M/2                   # center / group delay (samples)

# ideal low-pass sinc, centered and shifted to be causal
x = wc*(n - a)
h_ideal = np.where(x == 0, wc/np.pi, np.sin(x)/(np.pi*(n - a)))

# Hamming window, then taper the ideal taps
w = 0.54 - 0.46*np.cos(2*np.pi*n/M)
h = h_ideal * w
h /= h.sum()                 # normalize DC gain to exactly 1

# frequency response, to check against the spec
W   = np.linspace(0, np.pi, 512)
H   = np.array([np.sum(h*np.exp(-1j*om*n)) for om in W])
mag_dB = 20*np.log10(np.abs(H) + 1e-12)
\end{lstlisting}

\begin{figure}[H]
\centering
\begin{tikzpicture}
\begin{axis}[
  width=13.5cm, height=5.6cm,
  xmin=0, xmax=500, ymin=-80, ymax=8,
  xtick={0,100,150,200,300,400,500},
  ytick={0,-3,-40,-60}, yticklabels={$0$,$-3$,$-40$,$-60$},
  xlabel={frequency (Hz)}, ylabel={$|H|$ (dB)},
  axis lines=box, enlargelimits=false,
  title={\small Magnitude response of the 35-tap Hamming low-pass},
]
% passband corridor
\addplot[draw=none, fill=sigblue!12, domain=0:100, samples=2] {8} \closedcycle;
% stopband requirement line
\draw[specred, dashed, thick] (axis cs:200,-40)--(axis cs:500,-40);
\draw[gray, dashed] (axis cs:100,-80)--(axis cs:100,8);
\draw[gray, dashed] (axis cs:200,-80)--(axis cs:200,8);
% modeled response: smooth low-pass with stopband floor near -53 dB and lobes
\addplot[sigblue, very thick, samples=160, domain=2:500]
  { max( 20*log10( 1/sqrt(1+(x/162)^14) ) ,
         -53 + 6*abs(cos(deg(0.055*x))) ) };
\node[specred, anchor=south west] at (axis cs:210,-38) {\scriptsize $-40$\,dB spec};
\node[gray] at (axis cs:150,5) {\scriptsize $f_c$};
\end{axis}
\end{tikzpicture}
\caption{The designed filter (modeled curve). It is flat through the $100$\,Hz
passband, falls through the transition, and sits well below the $-40$\,dB spec line
(dashed red) by $200$\,Hz, with a stopband floor near $-53$\,dB set by the Hamming
window. Spec met.}
\label{ch16:fig-fir-resp}
\end{figure}

\begin{figure}[H]
\centering
\begin{tikzpicture}
\begin{axis}[
  width=13.5cm, height=4.6cm,
  xmin=0, xmax=1, ymin=-1.6, ymax=1.6,
  xtick={0,0.5,1}, xticklabels={$0$,$0.5$,$1.0$},
  ytick={-1,0,1},
  xlabel={time (s)}, ylabel={accel (a.u.)},
  axis lines=box, enlargelimits=false,
  legend style={font=\scriptsize, at={(0.98,0.97)}, anchor=north east},
  title={\small The low-pass at work: a slow tilt buried in fast jitter},
]
% noisy input: slow sine + high-frequency wiggle
\addplot[sigblue!45, samples=200, domain=0:1]
  {sin(deg(2*pi*x)) + 0.45*sin(deg(2*pi*40*x)) + 0.3*cos(deg(2*pi*55*x))};
% smoothed output: just the slow component, slightly delayed
\addplot[specred, very thick, samples=120, domain=0:1]
  {sin(deg(2*pi*(x-0.017)))};
\legend{raw input, after low-pass}
\end{axis}
\end{tikzpicture}
\caption{What the 35-tap low-pass does to a signal shaped like the accelerometer
problem: a slow $1$\,Hz tilt (the part we want) drowning in $40$--$55$\,Hz jitter.
The filter passes the slow component nearly untouched and erases the fast jitter,
delaying everything by the constant $17$-sample group delay. This is gravity
separated from vibration.}
\label{ch16:fig-fir-action}
\end{figure}

\begin{walkbox}
Suppose the verification plot showed the stopband only reaching $-35$\,dB --- one
nudge short of spec. What is the cheapest fix? Not more taps (Hamming's floor is
fixed by the \emph{window}, not the length; more taps only narrows the
transition). The right move is to switch to a window with deeper side lobes ---
Blackman ($-58$\,dB floor) --- and then, since Blackman's transition is wider,
bump the length back up to keep the $100$--$200$\,Hz transition. This is the window
method's whole logic in one decision: \emph{window sets the floor, length sets the
transition.}
\end{walkbox}

\begin{trapbox}
A subtle bookkeeping bug: after windowing, the DC gain $\sum_k h[k]$ is usually
close to $1$ but not exactly $1$, so a constant input comes out slightly scaled.
The one-line fix is to normalize, \texttt{h /= h.sum()} (for a low-pass). Forget
it and your gravity estimate on the accelerometer drifts by a percent or two ---
small, but it accumulates in an orientation estimate.
\end{trapbox}

\section{IIR design: borrow from the analog masters}
\label{ch16:iir}

FIR is built directly in the digital world. IIR design takes a different,
historically older route: \emph{steal a great analog filter and convert it to
digital.} Analog engineers spent a century perfecting smooth, efficient
frequency-response shapes out of resistors, capacitors, and inductors. Those
shapes are described by their \emph{poles and zeros} in the $s$-plane (the Laplace
domain of the previous chapter), and we can transplant them into the digital
$z$-plane.

\subsection{The Butterworth prototype: maximally flat}
\label{ch16:butterworth}

The most beloved analog prototype is the \textbf{Butterworth} low-pass. Its
defining property is geometric and easy to picture: it is \emph{maximally flat} in
the passband. Its magnitude-squared response is
\[
|H(\jj\Omega)|^2 \;=\; \frac{1}{1 + (\Omega/\Omega_c)^{2N}},
\]
where $N$ is the filter order and $\Omega_c$ the cutoff. Read that formula
geometrically. At $\Omega = 0$ the response is exactly $1$, and the first $2N-1$
derivatives are zero --- the curve leaves DC as flat as a table. At $\Omega =
\Omega_c$ it is always $1/\sqrt2$ ($-3$\,dB), regardless of $N$. Far past cutoff it
rolls off at $20N$\,dB per decade --- doubling $N$ doubles the steepness.

\begin{figure}[H]
\centering
\begin{tikzpicture}
\begin{axis}[
  width=13.5cm, height=5.0cm,
  xmin=0, xmax=3, ymin=-0.03, ymax=1.08,
  xtick={1}, xticklabels={$\Omega_c$},
  ytick={0,0.707,1}, yticklabels={$0$,$\tfrac{1}{\sqrt2}$,$1$},
  xlabel={frequency $\Omega/\Omega_c$}, ylabel={$|H|$},
  axis lines=box, enlargelimits=false,
  legend style={font=\scriptsize, at={(0.98,0.95)}, anchor=north east},
  title={\small Butterworth low-pass: higher order = sharper, still flat},
]
\addplot[sigblue, thick, samples=120, domain=0:3] {1/sqrt(1+(x)^(2*1))};
\addplot[specred, thick, samples=120, domain=0:3] {1/sqrt(1+(x)^(2*2))};
\addplot[quotientorange, thick, samples=120, domain=0:3] {1/sqrt(1+(x)^(2*4))};
\addplot[structurepurple, thick, samples=140, domain=0:3] {1/sqrt(1+(x)^(2*8))};
\draw[gray,dashed] (axis cs:1,-0.03)--(axis cs:1,1.08);
\draw[gray,dashed] (axis cs:0,0.707)--(axis cs:3,0.707);
\legend{$N=1$,$N=2$,$N=4$,$N=8$}
\end{axis}
\end{tikzpicture}
\caption{Butterworth magnitude for several orders. Every curve passes through
$-3$\,dB at $\Omega_c$ and is monotonic (no ripple). Higher order squares off the
corner. This single shape, with order chosen to meet a spec, covers a huge
fraction of real IIR designs.}
\label{ch16:fig-butter}
\end{figure}

\begin{intuitionbox}
The poles of an order-$N$ Butterworth sit on a semicircle in the left half of the
$s$-plane, perfectly evenly spaced (Figure~\ref{ch16:fig-butter-poles}). That geometric tidiness --- points equally
spread on a circle, exactly like the roots of unity from the DFT chapter --- is
\emph{why} the passband is so flat: the contributions of the poles balance to
cancel every low-order wiggle. Other prototypes move those poles onto an ellipse
(Chebyshev: steeper, but ripple appears) or add zeros (elliptic: steepest of all,
ripple in both bands). Butterworth is the ``no ripple anywhere'' choice.
\end{intuitionbox}

\begin{figure}[H]
\centering
\begin{tikzpicture}[scale=2.4]
\draw[->,gray!70] (-1.35,0)--(1.0,0) node[right]{\scriptsize $\sigma$};
\draw[->,gray!70] (0,-1.2)--(0,1.2) node[above]{\scriptsize $\jj\Omega$};
\draw[gray, thick] (0,0) circle (1);
% order-4 Butterworth poles: angles 112.5, 157.5, 202.5, 247.5 deg (left half)
\foreach \a in {112.5,157.5,202.5,247.5}{
  \node[structurepurple] at ({cos(\a)},{sin(\a)}) {\large $\times$};
}
\draw[gray!50,dashed] (0,0)--({cos(112.5)},{sin(112.5)});
\node[gray] at (-0.32,0.18) {\scriptsize $\tfrac{\pi}{8}$};
\node[structurepurple,anchor=west] at (0.15,-0.95) {\scriptsize order-4 poles};
\end{tikzpicture}
\caption{The four poles of an order-4 Butterworth low-pass in the analog
$s$-plane. They sit on a circle of radius $\Omega_c$, evenly spaced and confined
to the stable left half-plane. Equal spacing on the circle (the same geometry as
roots of unity) is what makes the passband maximally flat. Higher orders add more
poles around the same semicircle.}
\label{ch16:fig-butter-poles}
\end{figure}

\begin{historybox}
Stephen Butterworth, a British engineer, described this filter in a 1930 paper
titled simply ``On the theory of filter amplifiers.'' He wanted an amplifier whose
gain was as flat as possible across the passband for radio work, and he showed the
all-poles-on-a-circle design achieved it. He reportedly noted that the
mathematics, while elegant, required solving high-degree equations ``which can
only be done by hard work'' --- a reminder that before computers, even choosing a
filter order was labor. Today \texttt{scipy.signal.butter(N, Wc)} returns the
coefficients in microseconds.
\end{historybox}

\subsection{The bilinear transform: warping analog into digital}
\label{ch16:bilinear}

We have a gorgeous analog design $H_a(s)$. How do we make it digital, i.e. find an
$H(z)$ with the same shape? The bridge is the \textbf{bilinear transform}:
\[
s \;=\; \frac{2}{T}\,\frac{1 - z^{-1}}{1 + z^{-1}}.
\]
Substituting this into $H_a(s)$ turns the continuous transfer function into a
discrete one with the same general response. Why this particular substitution?
Recall from the last chapter that the analog $s$-plane's imaginary axis (real
frequencies) must map onto the digital $z$-plane's unit circle (digital
frequencies), and the stable left-half $s$-plane must map \emph{inside} the unit
circle (so stable analog filters become stable digital filters). The bilinear
transform is the simplest algebraic map that does exactly this.

\begin{happeningbox}
The bilinear transform takes the \emph{entire} infinite analog frequency axis
($\Omega$ from $0$ to $\infty$) and squashes it onto the \emph{finite} digital
range ($\omega$ from $0$ to $\pi$). Infinite into finite means the mapping must
compress as you go up: it is the warp
\[
\Omega \;=\; \frac{2}{T}\tan\!\Big(\frac{\omega}{2}\Big).
\]
Near $\omega = 0$ this is nearly linear ($\Omega \approx \omega/T$) --- digital and
analog frequencies agree. But as $\omega \to \pi$, the $\tan$ blows up, so high
digital frequencies get pulled from ever-higher analog frequencies. This
\emph{frequency warping} is the one quirk of the method: a Butterworth corner you
placed at a particular analog frequency lands at a slightly different digital
frequency than you would naively expect.
\end{happeningbox}

\begin{whybox}
The cure for warping is to ``pre-warp'': before designing the analog prototype,
deliberately distort the target cutoff by the \emph{inverse} of the warp, so that
after the bilinear transform squashes it, it lands exactly where you wanted. If
you want a digital cutoff at $\omega_c$, design the analog prototype for
$\Omega_c = \tfrac{2}{T}\tan(\omega_c/2)$. The warp and the pre-warp cancel at the
cutoff, and your $-3$\,dB point ends up precisely on target. Every library's
\texttt{butter} does this for you automatically; knowing it exists explains why the
band edges still come out right despite the nonlinear map.
\end{whybox}

\subsection{Worked design: a Butterworth low-pass in three lines}
\label{ch16:worked-iir}

Same spec target as before --- a low-pass on a $1000$\,Hz stream with cutoff near
$150$\,Hz --- but now IIR.

\begin{lstlisting}[language=Python]
from scipy.signal import butter, sosfilt
import numpy as np

fs = 1000.0
fc = 150.0                       # -3 dB cutoff (Hz)
N  = 4                           # order: ~24 dB/octave rolloff

# design: 'sos' = cascade of second-order sections (numerically robust)
sos = butter(N, fc/(fs/2), btype='low', output='sos')

# apply to a signal x
# y = sosfilt(sos, x)            # causal IIR filtering
\end{lstlisting}

Four coefficients-worth of work (two biquad sections) gives a rolloff that an FIR
would need dozens of taps to match. That is the IIR bargain: sharpness for almost
nothing. The bills are the two we already named --- nonlinear phase, and a
stability you must respect.

\begin{buildbox}
\textbf{Run IIR as biquads, never as one big polynomial.} A high-order IIR written
as a single long difference equation is a numerical minefield on a
fixed-point/16-bit micro: the coefficients span a huge dynamic range and tiny
rounding errors can push a pole outside the unit circle, turning your stable filter
into an oscillator. The fix is universal: factor the filter into a \emph{cascade of
second-order sections} (``biquads''), each a tiny $2$-pole/$2$-zero stage. SciPy's
\texttt{output='sos'} gives you exactly this. On an Arduino-class chip a biquad is
five multiplies and four adds per sample --- trivial --- and the cascade stays
stable. This is the standard way audio codecs, EQs, and our wrist controller's
EMG band-pass are actually run.
\end{buildbox}

\begin{trapbox}
\textbf{Stability is not optional with IIR.} An FIR is unconditionally stable; an
IIR is stable only if every pole sits strictly inside the unit circle. Quantizing
coefficients, or designing too high an order with too sharp a cutoff, can nudge a
pole onto or outside the circle, and the filter \emph{self-oscillates} or blows up
--- there is no input that fixes it. Always verify the designed poles' radii
($<1$), keep the order modest, and run as biquads. When in doubt, an FIR cannot do
this to you.
\end{trapbox}

\section{FIR versus IIR, decided by the spectrum plot}
\label{ch16:compare}

Put the two designs side by side and the choice becomes concrete. For the
\emph{same} $-3$\,dB cutoff, the order-$4$ Butterworth and the $35$-tap Hamming FIR
have comparable magnitude rolloff --- but their phase behavior is night and day.

\begin{figure}[H]
\centering
\begin{tikzpicture}
\begin{axis}[
  width=13.5cm, height=5.0cm,
  xmin=0, xmax=500, ymin=-2, ymax=30,
  xtick={0,100,150,200,300,400,500},
  ytick={0,17},
  xlabel={frequency (Hz)}, ylabel={group delay (samples)},
  axis lines=box, enlargelimits=false,
  legend style={font=\scriptsize, at={(0.98,0.95)}, anchor=north east},
  title={\small Group delay: FIR is flat (linear phase), IIR bends},
]
% FIR: constant group delay = M/2 = 17
\addplot[sigblue, very thick, samples=2, domain=0:500] {17};
% IIR: rising group delay near cutoff, then falling
\addplot[specred, thick, samples=120, domain=2:500]
  { 4 + 9/( (1+( (x-150)/70 )^2 ) ) };
\draw[gray,dashed] (axis cs:150,-2)--(axis cs:150,30);
\legend{FIR (35-tap Hamming), IIR (order-4 Butterworth)}
\end{axis}
\end{tikzpicture}
\caption{The deciding picture. The FIR delays every frequency by exactly $17$
samples (flat blue line) --- waveform shape preserved. The IIR's delay (red)
peaks near the cutoff: frequencies there lag the others, smearing shape. Same
magnitude job, very different effect on timing.}
\label{ch16:fig-groupdelay}
\end{figure}

\begin{aestheticbox}
This is why ``which filter?'' is never a one-size answer. If you only care
\emph{how loud} each frequency ends up --- audio tone shaping, an envelope, a power
estimate --- IIR's cheap sharpness wins and nobody hears the phase. If you care
\emph{when} things happen --- aligning sensor channels, preserving the shape of a
transient, measuring a delay --- FIR's flat group delay is worth every extra tap.
The spectrum plot tells you the magnitude story; the group-delay plot tells you the
timing story; a good engineer reads both before committing.
\end{aestheticbox}

\begin{applicationbox}
\textbf{The mains-hum notch.} A microphone or a biopotential amplifier picks up a
sharp $50$ or $60$\,Hz tone from the power grid. The ideal cure is a band-stop so
narrow it removes only the hum and leaves the music or muscle signal intact. This
is the textbook IIR job: a single biquad with a complex-conjugate pair of zeros
right on the unit circle at the hum frequency (forcing the gain to exactly zero
there) and a matching pole pair just inside the circle (so the notch is razor
thin), as in Figure~\ref{ch16:fig-notch-pz}. Two coefficient-pairs, a notch you
could not begin to match with a short FIR. The pole radius sets the notch width; push it toward $1$ for a narrower notch
--- but not \emph{onto} the circle, or you lose stability.
\end{applicationbox}

\begin{figure}[H]
\centering
\begin{tikzpicture}[scale=2.6]
% unit circle
\draw[gray, thick] (0,0) circle (1);
\draw[->,gray!70] (-1.25,0)--(1.3,0) node[right]{\scriptsize Re};
\draw[->,gray!70] (0,-1.25)--(0,1.3) node[above]{\scriptsize Im};
% notch angle ~ 60/500*pi rad ~ 21.6 deg for f0=60, fs=1000 (half-band=500)
% zeros on the circle, poles just inside at radius 0.92
\def\ang{38}
% zeros (on circle)
\node[specred] at ({cos(\ang)},{sin(\ang)}) {\large $\circ$};
\node[specred] at ({cos(-\ang)},{sin(-\ang)}) {\large $\circ$};
% poles (inside, radius 0.9)
\node[sigblue] at ({0.9*cos(\ang)},{0.9*sin(\ang)}) {\large $\times$};
\node[sigblue] at ({0.9*cos(-\ang)},{0.9*sin(-\ang)}) {\large $\times$};
% angle guide
\draw[gray!60,dashed] (0,0)--({cos(\ang)},{sin(\ang)});
\node[gray] at ({0.55*cos(\ang)+0.05},{0.55*sin(\ang)+0.12}) {\scriptsize $\omega_0$};
\node[specred,anchor=west] at (1.05,0.78) {\scriptsize $\circ$ zeros (on circle)};
\node[sigblue,anchor=west] at (1.05,0.62) {\scriptsize $\times$ poles (just inside)};
\end{tikzpicture}
\caption{Pole-zero plot of a mains-hum notch biquad. A conjugate pair of zeros
sits exactly \emph{on} the unit circle at the hum angle $\omega_0 = 2\pi f_0/f_s$,
forcing the gain to zero there. A conjugate pole pair sits at the same angle but
just inside the circle (radius near $0.9$); the closer the poles creep toward the
circle, the narrower the notch --- but a pole \emph{on} the circle would make the
filter unstable.}
\label{ch16:fig-notch-pz}
\end{figure}

\begin{goalbox}
\textbf{The wrist controller's filter stack, decided.}
\begin{itemize}
\item \emph{Accelerometer, separate gravity from motion}: a low-pass with cutoff
around $1$--$3$\,Hz isolates the slow gravity/tilt component; subtract it to get
the fast motion (Figure~\ref{ch16:fig-fir-action} shows this separation in
action). Channels must stay time-aligned across $x,y,z$ so a flick spike
looks simultaneous --- use a short \textbf{linear-phase FIR} so all three axes
share one constant delay.
\item \emph{EMG/FMG muscle band}: band-pass roughly $20$--$450$\,Hz to grab muscle
activity, then rectify and low-pass to get an envelope. Phase does not matter for
an envelope, so use cheap \textbf{IIR biquads}.
\item \emph{Mains hum}: a single \textbf{IIR notch} biquad at $50$ or $60$\,Hz.
\end{itemize}
And the limit, stated plainly: this stack cleans up and bands the signals,
but no amount of filtering lets a \emph{wrist-mounted IMU alone} recover individual
finger joint angles. Filtering separates gravity from motion and isolates frequency
bands; it cannot manufacture information the sensor never captured. Coarse,
discrete wrist and forearm gestures --- yes. Continuous per-finger pose needs
sensors on the fingers (rings) or richer forearm sensing (multi-channel EMG/FMG),
as we discuss in the hand-sensing chapter. A filter is a frequency-shaped scalpel,
not a source of new degrees of freedom.
\end{goalbox}

\section{Design checklist}
\label{ch16:checklist}

The whole chapter compresses to a procedure you can run every time.

\begin{enumerate}
\item \textbf{Write the spec}: band type; passband and stopband edges; allowed
ripple $\delta_p$; required attenuation $\delta_s$. Draw the tolerance corridor.
\item \textbf{Decide phase}: does waveform \emph{shape}/timing matter? Yes $\to$
lean FIR (linear phase). No, only magnitude $\to$ IIR is fine and cheaper.
\item \textbf{If FIR}: pick a window from your stopband attenuation (window sets
the floor), then a length from your transition width (length sets the steepness);
build the windowed sinc; normalize the DC gain.
\item \textbf{If IIR}: pick a prototype (Butterworth for no ripple), choose the
order from your rolloff/attenuation need, design via the bilinear transform (the
library pre-warps for you), and emit second-order sections.
\item \textbf{Verify}: plot the magnitude against the corridor \emph{and} the
group delay. For IIR, check every pole radius $< 1$.
\item \textbf{Implement}: FIR as a tapped sum; IIR as cascaded biquads. Confirm on
real data, on the real hardware, at the real sample rate.
\end{enumerate}

\begin{recoverybox}
If only one sentence survives: \emph{a filter is a frequency mask you design, FIR
shapes that mask with a windowed sinc and gives you exact linear phase at the cost
of taps, IIR borrows an analog prototype through the bilinear transform and gives
you cheap sharpness at the cost of phase and a stability you must watch.} Everything
else is detail in service of hitting a tolerance corridor.
\end{recoverybox}

\section{Exercises}
\label{ch16:exercises}

\begin{exercise}
A signal is sampled at $f_s = 8000$\,Hz. You want a low-pass that passes audio up
to $3000$\,Hz and strongly rejects everything above $3800$\,Hz. Give the passband
edge, the stopband edge, and the transition width, all in (a) Hz, (b) cycles per
sample, and (c) radians per sample. Why is $\omega = \pi$ the highest meaningful
frequency here, and what Hz does it correspond to?
\end{exercise}

\begin{exercise}
The ideal low-pass impulse response is $h[n] = \sin(\omega_c n)/(\pi n)$. (a) Use
l'Hopital or the small-angle limit to show $h[0] = \omega_c/\pi$. (b) Explain in
words why $h[n]$ being nonzero for all $n$ makes the ideal filter both non-causal
and infinitely long. (c) What single operation makes a truncated version causal,
and how many samples of delay does it cost?
\end{exercise}

\begin{exercise}
You truncate the ideal sinc to $21$ taps with a rectangular window and find the
first stopband ripple reaches $-21$\,dB. A colleague says ``just use $201$ taps,
the ripple will vanish.'' Explain why they are wrong (name the phenomenon), what
$201$ taps \emph{does} improve, and what actually lowers that first ripple.
\end{exercise}

\begin{exercise}
Show that a symmetric FIR, $h[k] = h[M-k]$, has linear phase. (Hint: group the
$k$ and $M-k$ terms in $H(\omega) = \sum_k h[k] e^{-\jj\omega k}$, factor out
$e^{-\jj\omega M/2}$, and show what remains is purely real.) State the constant
group delay in samples.
\end{exercise}

\begin{exercise}
Design, on paper, a $15$-tap Hamming-windowed low-pass with normalized cutoff
$\omega_c = 0.25\pi$. Write the formula for each tap $h[n]$ ($n=0,\dots,14$),
identify the center tap and its value, and state the filter's group delay. You do
not need to compute all the numbers --- show the setup and compute $h[7]$ (the
center) and $h[6]$ exactly.
\end{exercise}

\begin{exercise}
For a Butterworth low-pass, $|H(\jj\Omega)|^2 = 1/(1+(\Omega/\Omega_c)^{2N})$.
(a) Verify the gain at $\Omega = \Omega_c$ is $1/\sqrt2$ for every order $N$.
(b) Show the asymptotic rolloff far above cutoff is $20N$\,dB per decade of
frequency. (c) What order $N$ is needed to reach $-60$\,dB attenuation one decade
above the cutoff?
\end{exercise}

\begin{exercise}
The bilinear transform warps frequencies by $\Omega = \tfrac{2}{T}\tan(\omega/2)$.
For $T=1$, you want a digital cutoff at $\omega_c = 0.4\pi$. (a) What analog cutoff
$\Omega_c$ should you pre-warp to? (b) By what percentage does it differ from the
naive choice $\Omega_c = \omega_c$? (c) Explain why ignoring pre-warping matters
more for high cutoffs than low ones.
\end{exercise}

\begin{exercise}
You implement an order-$8$ Butterworth as a single $9$-coefficient difference
equation in $16$-bit fixed point and the output slowly grows without bound even
for a tiny input. (a) What has most likely happened to at least one pole? (b) Why
does a cascade of four biquads avoid this? (c) Could this failure ever happen to
an FIR filter? Justify.
\end{exercise}

\begin{exercise}
For the wrist controller, you must low-pass the three accelerometer axes to extract
gravity. A teammate proposes an order-$2$ IIR Butterworth (cheap) on each axis.
(a) What artifact will this introduce that could misalign the three axes during a
fast flick? (b) Propose an FIR alternative and state what you trade to get the fix.
(c) For the separate EMG envelope path, argue why the same teammate's IIR choice is
actually the right call there.
\end{exercise}

\begin{exercise}
Design a $60$\,Hz notch (band-stop) as a single biquad for $f_s = 1000$\,Hz.
(a) At what angle on the unit circle do you place the zeros, and why exactly on the
circle? (b) Where do you place the poles to make the notch narrow, and what
constrains how close to the circle they may go? (c) Sketch the resulting magnitude
response and mark the notch frequency and the $-3$\,dB notch width.
\end{exercise}

\begin{exercise}
(Open-ended build.) On an Arduino or Raspberry Pi, implement the $35$-tap FIR
low-pass from Section~\ref{ch16:worked-fir} as a circular buffer of the last $35$
samples. (a) Estimate the multiply-accumulate count per output sample. (b) At
$f_s = 1000$\,Hz, what fraction of the CPU does that consume on your board?
(c) Now implement the order-$4$ IIR as two biquads and compare the per-sample
cost. Which would you ship, and on which channel?
\end{exercise}
